\section{Kiểm chứng chéo (Cross-Dataset Validation)}
\label{sec:validation}

Để xác nhận tính tổng quát của phát hiện chính, chúng tôi chạy pipeline tối ưu (BGE-base + Adaptive
RRF) trên bốn bộ dữ liệu BEIR với các đặc điểm khác nhau. \textbf{nDCG@10 là metric chính}, paired
t-test với Shapiro-Wilk kiểm tra được sử dụng để kiểm định ý nghĩa thống kê. Ngoài ra, chúng tôi
fine-tune mô hình embedding riêng (BGE-small 33M) với RRF disagreement signal trên SciFact, đạt
\textbf{0.7909 nDCG@10} --- kết quả cao nhất trên SciFact, vượt cả BGE-base (109M) và vstash (0.7707).

\subsection{Đặc điểm các bộ dữ liệu}
\label{sec:validation:datasets}

\begin{table}[htbp]
\centering
\caption{Đặc điểm bốn bộ dữ liệu BEIR dùng trong kiểm chứng.}
\label{tab:dataset-characteristics}
\resizebox{\textwidth}{!}{%
\begin{tabular}{lrrrrr}
\toprule
Dataset & Docs & Test Q & Domain & Query style & BM25 nDCG@10 \\
\midrule
SciFact   & 5{,}183  & 300 & Scientific claims  & Formal, technical            & 0.6523 \\
NFCorpus  & 3{,}633  & 323 & Biomedical         & Information-dense abstracts  & 0.3079 \\
FiQA      & 57{,}638 & 648 & Financial QA       & Colloquial user questions    & 0.2167 \\
SciDocs   & 25{,}657 & 1{,}000 & Citation prediction & Paper titles as queries      & 0.1495 \\
\bottomrule
\end{tabular}
}
\end{table}

\subsection{Kết quả}
\label{sec:validation:results}

Bảng~\ref{tab:cross-dataset} tổng hợp kết quả pipeline trên bốn bộ dữ liệu với nDCG@10.
Bảng~\ref{tab:significance} (Phụ lục) trình bày đầy đủ các paired t-test. Bốn xu hướng
nhất quán: (1)~BGE-base luôn vượt trội SciNCL với biên độ lớn và có ý nghĩa thống kê (tất cả
$p<0.001$); (2)~Adaptive RRF không cải thiện có ý nghĩa so với BGE-base trên bất kỳ dataset nào
($p\ge0.139$); (3)~Cross-encoder không bao giờ cung cấp cải thiện có ý nghĩa, và làm giảm chất
lượng có ý nghĩa trên SciFact ($p=0.015$); (4)~Fine-tune BGE-small (33M) với RRF disagreement
signal trên SciFact đạt \textbf{0.7909 nDCG@10} --- vượt BGE-base (109M, +0.0533) và vượt vstash
reference (0.7707, +0.0202), cho thấy self-supervised fine-tuning là hướng đi hiệu quả để xây dựng
mô hình embedding riêng. Kết quả chi tiết cho cả 4 dataset được trình bày trong
Bảng~\ref{tab:cross-dataset}. Kết quả fine-tune trên SciFact được trình bày trong
Bảng~\ref{tab:pipeline-evolution} (Stage~9).

\subsection{Phân tích theo đặc điểm bộ dữ liệu}
\label{sec:validation:analysis}

\textbf{SciFact (formal claims).} BM25 mạnh nhất (nDCG@10 = 0.6523) vì thuật ngữ khoa học chính xác
khớp trực tiếp với abstract. BGE-base cải thiện đáng kể (+0.1737, $p<0.001$), và Adaptive RRF thêm
+0.0137 nhưng \emph{không} có ý nghĩa thống kê ($p=0.139$). Cross-encoder làm \emph{giảm} có ý
nghĩa ($-$0.0374, $p=0.015$).

\textbf{Fine-tune BGE-small (33M) trên SciFact} đạt nDCG@10 = \textbf{0.7909} --- kết quả cao nhất
trên SciFact trong tất cả pipeline. So với BGE-base (109M, 0.7376), cải thiện +0.0533; so với
vstash reference (0.7707), cải thiện +0.0202. MAP@10 = 0.7479, MRR@10 = 0.7564, Recall@10 = 0.9113.
Mô hình được huấn luyện từ 47K RRF disagreement triples (SciFact train, 5 epochs) và vượt mô hình
lớn hơn 3 lần (BGE-base, 109M). Kết quả này chứng minh self-supervised fine-tuning với tín hiệu
RRF disagreement có thể tạo mô hình embedding mạnh từ tập nhỏ và mô hình nhỏ.

\textbf{NFCorpus (biomedical).} BM25 yếu hơn (nDCG@10 = 0.3079) nhưng vẫn mang tín hiệu hữu ích.
BGE-base cải thiện +0.1444 ($p<0.001$) so với SciNCL. Adaptive RRF về cơ bản trùng với BGE-base
(+0.0024, $p=0.688$) --- hoàn toàn không có khác biệt. CE cũng không có hiệu quả ($p=0.752$).

\textbf{FiQA (colloquial QA).} Đây là bộ dữ liệu khó nhất cho BM25 (nDCG@10 = 0.2167) vì query dùng
ngôn ngữ đời thường, không khớp với văn bản tài chính trong corpus. BGE-base tận dụng semantic
matching để đạt 0.3909 --- cải thiện +0.3110 ($p<0.001$) so với SciNCL (0.0800). Adaptive RRF
giảm nhẹ ($-$0.0082, $p=0.484$) vì BM25 quá yếu, gần như tự động giảm trọng số BM25 về 0 nhưng
vẫn thêm một lượng nhỏ nhiễu.

\textbf{Kết luận:} Mô hình embedding hiện đại (BGE-base) là baseline an toàn và hiệu quả nhất khi
không có specialist riêng. Adaptive RRF là một cải thiện tùy chọn, không có ý nghĩa thống kê, hữu ích
nhất khi BM25 có tín hiệu mạnh. Cross-encoder không cần thiết và có thể gây hại trong pipeline mới.

\textbf{SciDocs (citation prediction).} Đây là bộ dữ liệu khác biệt nhất: citation prediction trên
25K bài báo khoa học. BM25 rất yếu (nDCG@10 = 0.1495) vì title queries không khớp từ vựng với full-text
papers. SciNCL đạt 0.1951, cao hơn BM25 (+30.5\%). BGE-base tiếp tục vượt trội SciNCL (+0.0196,
$p=0.008$, ${}^{**}$) nhưng biên độ nhỏ nhất trong 4 dataset, cho thấy trên citation prediction,
khoảng cách giữa embedding chuyên biệt khoa học (SciNCL) và embedding đa dụng (BGE) thu hẹp đáng kể.
Adaptive RRF giảm nhẹ ($-$0.0050, $p=0.371$, n.s.), và CE làm giảm có ý nghĩa
($-$0.0187, $p<0.001$, ${}^{***}$) --- xác nhận mô hình CE gây hại nhất quán. SciDocs củng cố
kết luận chính: BGE-base là lựa chọn mặc định hiệu quả khi không có specialist riêng; tầng bổ sung
chỉ đáng dùng nếu nó giải quyết trade-off specialist/generalist.

\textbf{Tổng kết kiểm chứng:} BGE-base là generalist an toàn trên transfer, trong khi fine-tune
BGE-small với RRF disagreement tạo ra SciFact specialist mạnh. Vì hai mục tiêu này xung đột, hướng
cuối cùng không phải fine-tune thêm một mô hình duy nhất, mà là Specialist-Generalist Adaptive Fusion:
giữ specialist trên batch giống SciFact và fallback sang BGE-base khi có bằng chứng distribution shift.
